\subsection{Teilaufgabe 2}
\subsubsection{Aufgabenstellung}
In der zweiten Teilaufgabe sollten wir ein Programm schreiben welches Sprunganweisungen in Java
 Sinnvoll verdeutlichen.

\subsubsection{Anforderungsdefinition}
\begin{enumerate}
	\item Verwendung von Sprunganweisungen.
	\item Mindestens ein switch-Anweisung.
\end{enumerate}

\subsubsection{Entwurf}
Im folgenden ist ein Entwurf für ein einfaches Login. Zuerst wird der Nutzer aufgefordert seine ID
anzugeben. Dabei wird geprüft ob es sich um eine Zahl handelt und ob diese Zahl grö\ss er gleich Eins ist.
Anschlie\ss end wird wird in einem switch überprüft ob eine solche ID existiert und ob das angegebene
Passwort auch das richtige ist. Wenn alles Stimmt wird eine Willkommens Nachricht  ausgegeben.
Andernfalls wird der Nutzer aufgefordert alles erneut einzutippen.
\begin{center}
	\includegraphics[width=0.8\textwidth]{uml/uml_c6_p2.pdf}
\end{center}

\subsubsection{Quelltext}
\paragraph{Sprunganweisungen.java}\
\lstinputlisting[language = Java , frame = trBL , escapeinside={(*@}{@*)}]{../chapter_06/src/main/java/chapter_06/Sprunganweisungen.java}

\subsubsection{Testdokumentation}
Bei Zahlen die kleiner als 1 waren, sowie Buchstaben, kam es zu einer entsprechenden Fehlermeldung. Anschließend konnte man den Wert erneut eintippen.

\subsubsection{Benutzungshinweise}
Nach dem aufrufen des Programms, wird der Nutzer aufgefordert seine NutzerID anzugeben,
sowie anschlie\ss end sein Passwort. Bei inkorrekter eingaben, wird man erneut aufgefordert
die Daten einzutippen.

\subsubsection{Anwendungsbeispiel}
Bei Erfolgreicher Anmeldung:
\begin{lstlisting}[frame = trBL , escapeinside={(*@}{@*)}]
[sebastian@laptop bin]$ java Sprunganweisungen 	
Wilkommen...!
ID      : 1
Passwort: hallo
Juhu, Sie haben sich eingeloggt
[sebastian@laptop bin]$ 
\end{lstlisting}
Bei inkorrekter Anmeldung:
\begin{lstlisting}[frame = trBL , escapeinside={(*@}{@*)}]
[sebastian@laptop bin]$ java Sprunganweisungen 	
Wilkommen...!
ID      : 12
Passwort: qwert
Ihre Angaben sind leider falsch, versuchen Sie es erneut.
ID      : -1
Passwort: hallo
Die Zahl muss größ er gleich 1 sein
ID      : 
[sebastian@laptop bin]$ 
\end{lstlisting}